\documentclass[11pt,a4paper]{article}
\usepackage[utf8]{inputenc}
\usepackage[T1]{fontenc}
\usepackage[romanian]{babel}
\usepackage[margin=2cm]{geometry}
\usepackage{array}
\usepackage{tabularx}
\usepackage{enumitem}
\usepackage{graphicx}
\usepackage{setspace}
\usepackage{xcolor}
\usepackage{hyperref}
\usepackage{amssymb}
\setlength{\parindent}{0pt}
\newcommand{\fillline}[1][6cm]{\underline{\hspace{#1}}}
\newcommand{\fillbox}[1][0.4cm]{\framebox[#1]{\rule{0pt}{#1}}}
\newcommand{\checkboxempty}{$\square$}

\begin{document}

\begin{tabular}{p{5cm} p{5cm} p{5cm}}
\textbf{PRIMĂRIA} & \textbf{*802012*} & Nr. \fillline[3cm] / \fillline[1.5cm] \\
\textbf{CLUJ-NAPOCA} & & \\
\end{tabular}

\vspace{0.5cm}

Către,

\begin{center}
\textbf{\underline{DIRECŢIA DE ASISTENŢĂ SOCIALĂ ŞI MEDICALĂ}}\\
\underline{Serviciul Asistenţa Persoanelor cu Dizabilităţi}
\end{center}

\vspace{0.3cm}

\begin{center}
\textbf{CERERE}\\
pentru eliberarea cardului-legitimaţie de parcare pentru persoanele cu handicap
\end{center}

\vspace{0.3cm}

\textbf{I.} (Se completează cu datele persoanei cu handicap.)

\vspace{0.2cm}

\hspace{0.5cm}Subsemnatul/Subsemnata:

\begin{enumerate}[leftmargin=1cm]
\item Numele şi prenumele \underline{ {{nume_titular}} }
\item CNP \underline{ {{cnp_titular}} }
\item Domiciliul: Cluj-Napoca, str. \fillline[7cm] nr. \fillline[1cm], bl. \fillline[1cm], sc. \fillline[1cm], et. \fillline[1cm], ap. \fillline[1cm], cod poştal \fillline[2cm]
\item Telefon \fillline[6cm]
\item E-mail \fillline[6cm]
\item Certificat de încadrare în grad de handicap (număr/serie/dată) \underline{ {{nr_certificat}} / {{valabilitate_certificat}} } -- grad \underline{ {{grad_handicap}} }
\end{enumerate}

\vspace{0.3cm}

\textbf{II.} (Se completează de către familie, asistentul personal, asistentul personal profesionist sau însoţitorul, pentru persoanele cu handicap grav sau accentuat, părinte, tutore, asistent maternal sau persoana care se ocupă de creşterea şi îngrijirea copilului cu handicap grav sau accentuat în baza unei măsuri de protecţie specială, stabilită în condiţiile legii.)

\begin{enumerate}[leftmargin=1cm]
\item Numele şi prenumele \fillline[10cm]
\item Domiciliul: localitatea \fillline[5cm], judeţ \fillline[4cm],\\
str. \fillline[6cm] nr. \fillline[1cm], bl. \fillline[1cm], sc. \fillline[1cm], et. \fillline[1cm], ap. \fillline[1cm], cod poştal \fillline[2cm]
\item Telefon \fillline[6cm]
\item E-mail \fillline[6cm]
\end{enumerate}

\vspace{0.2cm}

\hspace{0.5cm}Actul şi valabilitatea acestuia, prin care persoana este desemnată reprezentant legal, sau documentul care face dovada reprezentativităţii, conform pct. II \fillline[6cm] (Prezintă documentele de identitate în original).

\vspace{0.3cm}

\hspace{0.5cm}Solicit prin prezenta, în conformitate cu prevederile \underline{Legii nr. 448/2006} privind protecţia şi promovarea drepturilor persoanelor cu handicap, republicată, cu modificările şi completările ulterioare, şi cu prevederile \underline{Hotărârii Guvernului nr. 268/2007} pentru aprobarea \underline{Normelor metodologice} de aplicare a prevederilor \underline{Legii nr. 448/2006} privind protecţia şi promovarea drepturilor persoanelor cu handicap, cu modificările şi completările ulterioare, \textbf{eliberarea unui card-legitimaţie de parcare pentru persoanele cu handicap.}

\vspace{0.3cm}

\hspace{0.5cm}Declar pe propria răspundere, sub sancţiunea falsului în declaraţii prevăzut de \underline{Codul penal}, că nu am ridicat cardul-legitimaţie de parcare pentru persoanele cu handicap de la altă autoritate emitentă.

\vspace{0.3cm}

\hspace{0.5cm}Sunt de acord cu prelucrarea datelor cu caracter personal în conformitate cu legislaţia în vigoare.

\vspace{1cm}

Data: \fillline[3cm] \hfill Semnătura: \fillline[4cm]

\vfill

\hrule
\vspace{0.2cm}
www.primariaclujnapoca.ro\\
telefon: 0264/596.030\\
Cluj-Napoca, str. Moţilor, nr. 3 \hfill \textbf{*802012*}

\vspace{0.2cm}
\hrule
\vspace{0.2cm}
\footnotesize
Timp estimativ de completare: \textbf{5 minute}\\
Datele personale care vă sunt solicitate prin prezenta cerere vor fi prelucrate numai în vederea procesării și soluționării solicitării dumneavoastră. Primăria municipiului Cluj-Napoca garantează securitatea procesării datelor și arhivarea acestora în conformitate cu prevederile legale în vigoare. Responsabilul Primăriei municipiului Cluj-Napoca cu protecția datelor poate fi contactat pe adresa de email dpo@primariaclujnapoca.ro.\\
În conformitate cu Regulamentul nr. 679/2016 aveți dreptul de a solicita Primăriei municipiului Cluj-Napoca, în ceea ce priveşte datele cu caracter personal referitoare la persoana vizată, accesul la acestea, rectificarea sau ştergerea acestora sau restricţionarea prelucrării sau a dreptului de a vă opune prelucrării, precum şi a dreptului la portabilitatea datelor.

\end{document}
